% Figure for Electromagnetism §A12.1: a charging capacitor; the loop C around the wire bounds a flat surface S1 pierced by the current and a bag-shaped surface S2 passing between the plates, where there is no current but a changing electric field.
\begin{tikzpicture}[>={Stealth[length=2.2mm]}, font=\small]
  % surface S2 (bag around the left plate, closing between the plates)
  \fill[figG!12] (1.5,1.3) .. controls (2.6,1.75) and (3.3,1.8) .. (3.3,1.2) -- (3.3,-1.2) .. controls (3.3,-1.8) and (2.6,-1.75) .. (1.5,-1.3) -- cycle;
  \draw[figG, thick] (1.5,1.3) .. controls (2.6,1.75) and (3.3,1.8) .. (3.3,1.2) -- (3.3,-1.2) .. controls (3.3,-1.8) and (2.6,-1.75) .. (1.5,-1.3);
  % wire and plates
  \draw[very thick] (-0.6,0) -- (3.0,0);
  \draw[very thick] (3.6,0) -- (5.6,0);
  \fill[black!70] (2.95,-1.0) rectangle (3.05,1.0);
  \fill[black!70] (3.55,-1.0) rectangle (3.65,1.0);
  \node[font=\scriptsize] at (2.8,1.2) {$+Q$};
  \node[font=\scriptsize] at (3.85,1.2) {$-Q$};
  % E field between plates
  \foreach \y in {-0.7,-0.35,0.35,0.7} \draw[->, figR, thin] (3.08,\y) -- (3.52,\y);
  \node[figR, font=\scriptsize, anchor=west] at (3.9,-0.7) {$\vec E$ growing};
  % current arrows
  \draw[->, thick] (-0.4,0.25) -- (0.5,0.25) node[midway, above, font=\scriptsize] {$I$};
  \draw[->, thick] (4.4,0.25) -- (5.3,0.25) node[midway, above, font=\scriptsize] {$I$};
  % surface S1 (flat disk, seen almost edge-on) and loop C
  \fill[figB!25] (1.5,0) ellipse (0.32 and 1.3);
  \draw[figB, thick] (1.5,0) ellipse (0.32 and 1.3);
  \draw[very thick] (1.18,0) -- (1.82,0);
  \node[font=\scriptsize] at (1.0,1.35) {loop $C$};
  \node[figB, font=\scriptsize] at (1.5,-1.6) {$S_1$};
  \node[figG, font=\scriptsize] at (2.7,-1.75) {$S_2$};
\end{tikzpicture}
